\documentclass{article}
\usepackage{graphicx} % Required for inserting images

\title{AI Optical Biosensors for Non-invasive Neurotransmitter Detection }
\author{Divya Warule}
\date{April 2026}

\begin{document}

\maketitle

\section{Introduction}
Neurotransmitters are the chemical messengers inside the brain that do communication between the neurons. 
They also follow certain cycle , understanding this cycle is essential for making biosensor that will detect the chemicals .

\section{Recent advancements}
In recent years advancements have been done in the neurotransmitters like integration of ai. machine learning  and deep learning which helps to seperate overlapping signals , analyze the complex signals given by the biosensors and improve accuracy .

AI also helps real time monitoring and prediction .
After  that optical biosensors came . these sensors use light based technique like fiber optics . When neurotransmitter interact with the sensor surface , they make changes in the light properties like intensity , wavelength etc. this changes are converted into measurable signals .

After that came nanotechnology which played a key role in improving the sensor performance as the nanoparticles increase the surface areas and enhances the interaction with target molecules which increases the sensitivity and selectivity . they also helps in developing small and portable sensors .

Then comes the modern research which focuses on non-invasive techniques rather than invasive implants .(non invasive are the sensors on body surface). 
Nanotransmitters can be detected using the bodyfluids like blood , saliva or sweat.

\section{Limitations}
With advancments limitations also exist . Sensors perform well in laboratory conditions but fails in real world . Long tern stability and biocompatibility remains challenge . AI modles require  large and high quility dataset also often lack transparency . there in also lack of standardization .

\end{document}


